\documentclass[10pt]{article}
\usepackage[margin=20mm]{geometry}
\usepackage{amsmath,amssymb,fontspec,unicode-math,xurl,hyperref}
\IfFontExistsTF{DejaVu Serif}{\setmainfont{DejaVu Serif}}{\setmainfont{Latin Modern Roman}}
\setmathfont{latinmodern-math.otf}
\setlength{\parindent}{0pt}\setlength{\parskip}{6pt}\sloppy
\begin{document}
Abstract: The paper’s finite-penalty stochastic optimal controller and its training distribution do not describe the same stochastic process. A one-dimensional example gives true variance \(3/8\) where the stated marginal has variance \(1/4\); a second admissible example makes the claimed transition innovation variance negative. This invalidates a central derivation of the finite-penalty framework and warrants Major revision. The deterministic mean and the infinite-penalty bridge limit survive these checks. A separate printed reverse-update sign and scale inconsistency is Minor because the inspected implementation uses the consistent noise convention. Empirical image-quality results were not reproduced. This is an initial review; adjudication remains pending.\par
Author: AI-assisted technical review\par
Date:\par
Title: Technical Verification of UniDB\par
\# Paper and evidence scope\par
The reviewed work is Kaizhen Zhu, Mokai Pan, Yuexin Ma, Yanwei Fu, Jingyi Yu, Jingya Wang, and Ye Shi, UniDB: A Unified Diffusion Bridge Framework via Stochastic Optimal Control, ICML 2025, I25-43713 (0344) [Retained latex: \textbackslash{}cite\{paper\}]. All 28 proceedings-PDF pages, including appendices A.1–A.9, B and C, were read. All ten original figures and relevant equation/table pages were visually inspected. Page numbers below refer to one-based PDF pages. The original SHA-256 is\par
.\par
The author repository was inspected at commit (14 August 2026). Static reading covered the complete SDE utility, denoising model, matching loss, super-resolution train/test scripts, corresponding differences for deraining/inpainting, selected configuration files, paired-data loading, image conversion and PSNR/SSIM routines. The README supplies a GOU implementation; no released VE/VP implementation was inspected. The acquired source packet contains 45 files, all verified against its manifest. No author code, model, data download, training or inference was run. Independent computations use only reviewer-written scalar formulas and symbolic identities.\par
\# What survives the mathematical checks\par
Write \(y\) for the prescribed target, \(F_t=\int_0^t f_s\,ds\), \(G_t^2=\int_0^t e^{-2F_s}g_s^2\,ds\), and \[d_t=\gamma^{-1}+e^{2F_T}(G_T^2-G_t^2).\] For deterministic scalar coefficients, positive finite \(\gamma\), and the additive linear diffusion in (11), the feedback in (13) is the usual linear-quadratic control. Its deterministic trajectory (14) also gives the controlled stochastic mean. Certainty equivalence justifies this feedback relation; it does not remove the covariance generated by the Brownian forcing.\par
As \(\gamma\to\infty\) and for \(t<T\) with nondegenerate remaining noise, the Gaussian transition-density gradient gives the controller in Theorem 4.2. The reviewed counterexamples below do not refute that limit. Proposition 4.3 concerns the deterministic objective (12). A lower optimized objective with a softer terminal penalty is compatible with nonzero terminal mismatch. For example, \(f=h=0\), \(g=T=1\), \(x_0=0\), \(y=1\) gives constant optimal deterministic control \(u=\gamma/(1+\gamma)\) and cost \(\gamma/[2(1+\gamma)]\), compared with hard-constraint energy \(1/2\). This cost comparison by itself does not prove superior perceptual image quality.\par
\# Major issue: finite-penalty marginals and transitions\par
Problem 1 (The training variance is incompatible with the controlled SDE).  Section 4.3, equation (18), and appendix A.4 (pages 5 and 17) combine the finite-\(\gamma\) mean with a hard-bridge variance that does not depend on \(\gamma\).\par
Equation reference: case, actual. Take a scalar instance satisfying the stated linear, additive-noise and positive-penalty setting: \[
 f=h=0,\qquad g=1,\quad T=\gamma=1,\quad x_0=0,\quad y=1.\] Substitution into (13) and (11) yields \[dX_t=\frac{1-X_t}{2-t}\,dt+dW_t,\qquad X_0=0.\] Its mean and variance satisfy \[\mu'_t=\frac{1-\mu_t}{2-t},\qquad
 v'_t=-\frac{2v_t}{2-t}+1,\qquad \mu_0=v_0=0,\] and hence \[
 \mu_t=t/2,\qquad v_t=t(2-t)/2.\] Equation (18) has this same mean, but \(v_t^{\rm paper}=t(1-t)\). At \(t=1/2\) these variances are \(3/8\) and \(1/4\), respectively. At \(T=1\) the true variance is \(1/2\) while the printed variance is zero. The disagreement therefore exists in the interior and is not solely a convention at a singular endpoint. Symbolic substitution verifies both moment ODEs in [actual] (\#actual); substituting the printed variance into the variance ODE gives a nonzero residual.\par
Equation reference: cov. More generally, integrating the fundamental solution \(e^{F_t}d_t/d_0\) gives the actual conditional-on-\(x_0,y\) covariance \[
 \operatorname{Var}(X_t\mid x_0,y)
 =e^{2F_t}G_t^2\frac{d_t}{d_0}\,I.\] The target \(y\) here is a fixed parameter of the feedback, not an additional hard conditioning event \(X_T=y\). At infinite penalty [cov] (\#cov) reduces to the hard-bridge expression. At finite penalty it generally does not.\par
Appendix A.4 applies Bayes’ formula to the uncontrolled Gaussian transitions (57). That computes a process conditioned to hit its endpoint exactly. Its mean must then also be the hard-conditioned mean. Keeping the soft-control mean while importing the hard-conditioned variance is not a derivation of the law of (11)/(13). The paper explicitly invokes stochastic interpolants to introduce a vanishing-endpoint noise term. Such an interpolant can be defined in its own right, but its compatible drift and joint transitions must be derived; it is not automatically the controlled SDE or its reverse (15).\par
\#\# The claimed transition can have negative variance\par
Appendix A.5, equation (63), further assumes that the marginal construction admits a Gaussian Markov transition with innovation variance \[q_{s,t}=v_t^{\rm paper}-v_s^{\rm paper}(a_t/a_s)^2,
 \qquad a_t=e^{F_t}d_t.\] Choose \(f=h=0\), \(g=1\), \(T=10\), \(\gamma=1/10\), \(s=8\) and \(t=9\). These are ordinary interior, unit-spaced times with a strictly positive finite penalty. Here \[a_t/a_s=11/12,\quad v_s^{\rm paper}=8/5,\quad
 v_t^{\rm paper}=9/10,
 \qquad q_{s,t}=\frac9{10}-\frac85\left(\frac{11}{12}\right)^2=-\frac49.\] No real Gaussian innovation has this variance. Thus interpreting (18) as an arbitrary interpolant alone does not rescue the specific joint construction and posterior mean used in (19). The issue is a failed central training derivation, not merely a small error in a terminal approximation.\par
\#\# Terminal-distribution interpretation and scope\par
For the GOU specialization, let \(S=\lambda^2(1-e^{-2\bar\theta_T})\). The finite-control terminal variance is \(S/(1+\gamma S)\). Proposition 4.5, equation (23), accounts only for squared mean displacement. Under \(\theta=\lambda=\gamma=T=1\) and \(x_0=y=0\), its right-hand side is zero, whereas the actual stochastic terminal mean-square error is \[\mathbb E|X_T-y|^2=\frac{1-e^{-2}}{2-e^{-2}}\simeq0.46371.\] For this point target the same quantity is the squared \(W_2\) distance. Equation (23) remains a valid statement about the deterministic mean. Its use in Figure 4 as a full controlled-terminal distribution distance requires correction. This is a downstream manifestation of the same covariance problem, not a separate Major finding.\par
The released uses a \(\gamma\)-independent standard deviation at lines 143–146 and 227–228, a \(\gamma\)-dependent mean at 190–194, the innovation square root at 202–203, and these together to generate training states at 363–375. The code therefore corroborates the hybrid construction in the inspected version. However, training samples \(t\in\{1,\ldots,99\}\), and the supplied configurations use \(\gamma^{-1}=10^{-7}\) or \(10^{-8}\). An independent evaluation of their scalar schedule gives strictly positive innovations at those sampled times (minimum about \(4.88\times10^{-6}\)). The negative endpoint innovation is excluded from training; sampling also special-cases \(t=100\). Consequently, the broad counterexample is not evidence that these default runs produced NaNs or that reported image scores are wrong. No historical-run correspondence is established.\par
The assessment is Major: the connection from finite-penalty SOC to the training process is a central contribution, and the claimed transition derivation fails under its own admissible parameters. Repair requires choosing a consistent process: use [cov] (\#cov) and rederive transition targets, endpoint initialization and reverse sampling, or explicitly develop a different interpolant with a compatible drift and joint law. Either repair needs new analysis; empirical performance of a repaired implementation remains untested.\par
\# Minor issue: printed reverse-noise convention\par
Problem 2 (Equations and pseudocode disagree with the declared score parametrization).  Equation (19) and appendix A.5 use a positive noise-prediction term after specifying the score as \(-\epsilon_\theta/\bar\sigma'_t\). Algorithms 1 and 2 additionally divide by \((\bar\sigma'_t)^2\).\par
Let \(b_t=f_tx_t+h_tm+g_tu^*_{t,\gamma}\). Equation (15) and a backward time step \(\Delta>0\) give \[\mathbb E[X_{t-\Delta}\mid X_t=x]
 \simeq x-\Delta b_t+\Delta g_t^2\nabla_x\log p_t(x)
 =x-\Delta b_t-\Delta\frac{g_t^2}{\bar\sigma'_t}\epsilon_\theta.\] At unit step, the printed equation (19) instead adds \(g_t^2\epsilon_\theta/\bar\sigma'_t\). Equation (66), page 19, exhibits the sign change when substituting the declared score. Algorithms 1 and 2, page 6, visibly use \(+g_t^2\epsilon_\theta/(\bar\sigma'_t)^2\). For \(g=1\), \(\bar\sigma'=1/2\), \(\epsilon_\theta=1\), the score contributions are respectively \(-2\), \(+2\), and \(+4\). A reversed definition of predicted noise could change a sign, but it would contradict the explicit score parametrization and would not by itself repair the extra denominator power.\par
The released implementation is important counterevidence. Lines 274–276 set ; lines 254–269 subtract \(g^2\) times that score inside the drift; lines 37–38 subtract that drift. Sampling follows the same sign and standard-deviation convention (322–345). The denoising model calls these functions at 134–140. Hence this verified defect is Minor, confined to the printed equations/pseudocode relative to their declared convention. Correct the sign, divide by the standard deviation, and state the time-step factors consistently. It is not a demonstrated wrong-sign implementation or an explanation of historical results.\par
\# Empirical evidence and remaining limits\par
Tables 1–6, the hyperparameter study, implementation details and all qualitative panels were examined. Table 1’s super-resolution SDE comparison improves LPIPS from 0.220 to 0.179 and FID from 20.85 to 16.21 while reducing PSNR from 26.89 to 25.46 and SSIM from 0.7478 to 0.6856. The paper explicitly prioritizes perceptual scores in section 5.3. These data support a reported tradeoff rather than uniform metric dominance. The observed panels show the claimed local texture/sharpness differences, but visual inspection cannot certify the benchmark values or the causal SOC explanation.\par
The released configurations specify 900,000 iterations while appendix B says 1.2 million, and the shared matching loss is unweighted L1. These release-versus-description differences are recorded as reproducibility limits, not evidence that the published experiments used the wrong budget or objective. Training logs, checkpoints, raw results, FID feature computation and sampling variability were not available in the reviewed packet. Network internals and data utilities beyond the listed paths were not exhaustively audited. The review does not establish statistical significance or reproduce any image metric.\par
The two confirmed findings require different responses: a coherent replacement derivation for the finite-penalty stochastic process, and local corrections to the printed reverse-update convention. The hard-bridge limit, the deterministic mean calculation, and the reported empirical measurements are not refuted by the supplied counterexamples. Final scientific adjudication and independent acceptance remain pending.\par
9 K. Zhu et al., UniDB: A Unified Diffusion Bridge Framework via Stochastic Optimal Control, ICML 2025. https://proceedings.mlr.press/v267/zhu25o.html. UniDB authors, UniDB repository, inspected commit . https://github.com/UniDB-SOC/UniDB.\par
\end{document}
